% Vietnamese translation for OI Public Library
% Source-File: IOI中国国家候选队论文/2008/Day1/8.吕子鉷《浅谈最短径路问题中的分层思想》/浅谈最短径路问题中的分层思想.ppt
\documentclass[11pt]{article}
\usepackage{oipl-vn}

\newcommand{\Oh}{\mathcal{O}}

\title{Bản trình chiếu: Phân tầng trong bài toán đường đi ngắn nhất}
\author{Lu Zihong}
\date{2008}

\begin{document}
\maketitle

\origtitle{Layered thinking in shortest paths}
\sourcepath{IOI China candidate team papers / 2008 / Day 1 / Lu Zihong.ppt}

\section{Phạm vi bản dịch}

Các slide đi cùng bài viết về tư tưởng phân tầng trong bài toán đường đi ngắn
nhất. Tệp PPT gốc có nhiều hình vẽ và mốc độ phức tạp, nhưng phần chữ trích
được cho thấy cấu trúc chính gồm năm ví dụ: \textit{Rescue}, \textit{Fence},
\textit{Cow Relay}, \textit{Bicriterial Routing} và \textit{Roads}. Bản ghi
chú này dịch và hệ thống hóa nội dung theo đúng mạch đó.

Ý tưởng chung là sao chép một đồ thị thành nhiều tầng. Mỗi tầng giữ một phần
trạng thái bổ sung: tập chìa khóa, tính chẵn lẻ của số lần cắt tia, số đường
đã loại bỏ, tổng phí, hoặc tổng chiều dài. Khi đại lượng ấy được đưa vào đỉnh,
bài toán có thể trở lại thành đường đi ngắn nhất hoặc quy hoạch động trên đồ
thị phân tầng.

\section{Rescue: tầng là tập chìa khóa}

Bài \textit{Rescue} có mê cung \(N\times M\). Giữa hai ô kề nhau có thể không
có vật cản, có tường, hoặc có cửa khóa. Có \(P\) loại chìa khóa, \(P\le 10\);
một số ô chứa chìa. Cần đi từ \((1,1)\) tới \((N,M)\) trong thời gian ngắn
nhất.

Nếu không có cửa khóa, mỗi ô là một đỉnh và mỗi bước đi là một cạnh trọng số
\(1\). Cửa khóa làm cho khả năng đi qua cạnh phụ thuộc vào tập chìa đã nhặt.
Vì vậy ta tạo \(2^P\) tầng, mỗi tầng ứng với một tập chìa khóa.

Trong một tầng, ta chỉ giữ những cạnh cửa mà tập chìa của tầng đó mở được.
Nếu ô hiện tại có chìa loại \(q\), thêm cạnh trọng số \(0\) từ bản sao của ô
trong tầng hiện tại sang bản sao trong tầng đã thêm chìa \(q\). Đồ thị mới có
\(\Oh(2^P NM)\) trạng thái. Vì trọng số chỉ là \(0\) và \(1\), có thể dùng
BFS \(0/1\) hoặc thuật toán đường đi ngắn nhất tương đương.

Điểm chính của slide là cách biến ``đã nhặt những chìa nào'' từ lịch sử của
đường đi thành một phần của đỉnh hiện tại.

\section{Fence: tầng là tính chẵn lẻ}

\textit{Fence} yêu cầu tìm hàng rào ngắn nhất bao quanh tòa nhà chính nhưng
không đi vào vùng cấm quanh các công trình khác. Các vùng cấm là hình chữ nhật
song song trục, và hàng rào cũng gồm các đoạn song song trục.

Trước hết dựng một đồ thị hình học: đỉnh là các góc hình chữ nhật; nối hai
đỉnh bằng cạnh nếu đoạn thẳng song song trục giữa chúng không đi xuyên qua
vùng cấm. Trọng số cạnh là độ dài đoạn. Một chu trình trong đồ thị có thể là
một hàng rào, nhưng ta còn phải biết nó có bao quanh tòa nhà chính hay không.

Slide dùng tiêu chuẩn chẵn lẻ: kẻ một tia từ điểm cần bao quanh; một đường
khép kín bao quanh điểm đó khi số lần cắt tia là lẻ. Do đó ta tách đồ thị
thành hai tầng. Cạnh nào cắt tia sẽ nối sang tầng còn lại; cạnh nào không cắt
tia sẽ nằm trong cùng tầng.

Khi đó một chu trình bao quanh điểm tương đương với một đường đi từ một bản
sao của đỉnh \(v\) tới bản sao \(v\) ở tầng còn lại. Chạy đường đi ngắn nhất
và lấy nhỏ nhất trên mọi \(v\) cho ta độ dài hàng rào tối ưu. Đây là ví dụ
điển hình nơi tầng không biểu diễn số lần đầy đủ, mà chỉ biểu diễn phần cần
thiết của nó: chẵn hay lẻ.

\section{Cow Relay: tầng là tập đường đi còn lại}

Trong \textit{Cow Relay}, cần tìm đường đi thứ \(k\) ngắn nhất giữa hai đỉnh,
với \(k\le 40\). Cách dùng hàng đợi ưu tiên để mở rộng các đường đi ứng viên
có thể giải bài, nhưng slide trình bày một cách nhìn phân tầng để loại dần
đường đi đã chọn.

Thêm nguồn \(s\) nối tới đỉnh đầu và đỉnh cuối nối tới \(t\). Sau khi tìm
đường ngắn nhất \(p_1\) từ \(s\) tới \(t\), ta tạo một tầng sao chép. Các cạnh
không thuộc \(p_1\) được phép chuyển từ tầng cũ sang tầng mới; các cạnh thuộc
\(p_1\) thì không. Vì vậy mọi đường từ \(s\) tới \(t'\) trong tầng mới tương
ứng với một đường cũ nhưng không thể trùng đúng \(p_1\). Đường ngắn nhất mới
chính là đường ngắn thứ hai.

Lặp ý tưởng này cho các đường đã tìm được sẽ thu được đường thứ \(k\). Bản
trình chiếu ghi các mốc tối ưu hóa: không cần sao chép toàn bộ đồ thị ở mọi
vòng, chỉ cần sao chép phần có thể tạo ra nhánh khác; cây đường đi ngắn nhất
của phần cũ được giữ lại. Với cách triển khai đó, độ phức tạp hướng tới
\[
  \Oh(km),
\]
và bộ nhớ khoảng \(\Oh(kn+m)\).

Thông điệp ở đây hơi khác hai ví dụ trước: tầng không phải trạng thái vật lý
của người đi, mà là trạng thái của không gian lời giải sau khi đã loại một số
đường.

\section{Bicriterial Routing: tầng là chi phí}

\textit{Bicriterial Routing} xét mạng đường có hai trọng số trên mỗi cạnh:
thời gian và phí. Một tuyến bị loại nếu có tuyến khác vừa không chậm hơn vừa
không đắt hơn, và tốt hơn ở ít nhất một tiêu chí. Cần đếm các cặp
\((\text{phí},\text{thời gian})\) không bị trội.

Ta chọn phí làm chỉ số tầng. Nếu tổng phí hiện tại là \(q\), một cạnh có phí
\(w\) và thời gian \(t\) nối từ tầng \(q\) sang tầng \(q+w\), với trọng số
đường đi là \(t\). Như vậy, tại mỗi tầng phí, khoảng cách ngắn nhất tới đích
cho ta thời gian tốt nhất ứng với mức phí ấy.

Cách trực tiếp là chạy Dijkstra trên toàn đồ thị mở rộng. Nhưng đồ thị có
cấu trúc một chiều theo phí: cạnh không đi từ tầng cao xuống tầng thấp. Vì
vậy có thể xử lý các tầng theo thứ tự phí tăng dần và chỉ giải bài toán nhỏ
trong từng tầng. Mốc trong bài viết là giảm từ dạng
\[
  \Oh(ncm\log(n^2c))
\]
xuống
\[
  \Oh(ncm\log n),
\]
trong đó \(c\) là chặn phí cạnh.

Sau khi có thời gian tốt nhất theo từng mức phí, chỉ cần quét các cặp để bỏ
những cặp bị trội.

\section{Roads: đổi trục phân tầng}

\textit{Roads} cũng có hai đại lượng trên mỗi cạnh: chiều dài và phí. Bob có
ngân sách không quá \(k\) và muốn tới đích với chiều dài nhỏ nhất. Mô hình
đầu tiên giống ví dụ trước: dùng tổng phí làm tầng, chạy Dijkstra trên đồ thị
mở rộng. Slide ghi mốc
\[
  \Oh(k(kn^2+m))
\]
cho cách thô và
\[
  \Oh(k(n^2+m))
\]
khi xử lý tầng phí theo thứ tự.

Tuy nhiên, bài này có chặn nhỏ trên chiều dài cạnh. Vì vậy slide đổi trục:
dùng tổng chiều dài làm tầng, còn giá trị lưu trong trạng thái là phí nhỏ
nhất. Đặt
\[
  f[i,j]=\text{phí nhỏ nhất để tới thành phố }j
  \text{ với tổng chiều dài đúng }i.
\]
Điều kiện đầu là \(f[0,1]=0\), các trạng thái khác là vô hạn. Với cạnh
\((j_0,j)\) có chiều dài \(\ell\) và phí \(w\), cập nhật
\[
  f[i,j]=\min\{f[i,j],\,f[i-\ell,j_0]+w\}.
\]
Đáp án là \(i\) nhỏ nhất sao cho \(f[i,N]\le k\). Nếu \(L\) là chặn chiều dài
cạnh, độ phức tạp trở thành
\[
  \Oh(nLm).
\]

Ví dụ này cho thấy không phải đại lượng nào cũng nên làm tầng. Chọn phí làm
tầng là tự nhiên, nhưng chọn chiều dài làm tầng mới khai thác tốt hơn giới
hạn dữ liệu.

\section{Kết luận}

Phân tầng có hai vai trò trong các slide. Thứ nhất, nó giúp dựng mô hình:
tập chìa khóa, tính chẵn lẻ, hoặc tập đường đi còn lại được biến thành một
phần của đỉnh. Thứ hai, nó giúp tối ưu: khi cạnh chỉ đi cùng tầng hoặc sang
tầng cao hơn, ta có thứ tự xử lý tự nhiên và không cần xem đồ thị mở rộng như
một khối đen.

Khi áp dụng, cần trả lời ba câu hỏi: đại lượng nào gây khó cho đồ thị ban
đầu; đại lượng ấy có bao nhiêu giá trị đủ nhỏ để làm tầng; và cấu trúc giữa
các tầng có thể khai thác thêm hay không. Nếu chỉ sao chép đồ thị mà không
khai thác cấu trúc, phân tầng có thể làm bài toán lớn hơn; nếu chọn đúng
trục, nó biến ràng buộc lịch sử thành một mô hình đường đi ngắn nhất rõ ràng.

\end{document}
